% ============================================================================
% Fichier  : partie1/P01_C03_ethique_neutralite.tex : FORN (LPI342)
% Sources : ISO/IEC 27037 (impartialité) ; ACPO ; déontologie de l'expertise.
% ============================================================================

\chapter{Éthique, Neutralité, Impartialité}
\label{chap:ethique}
\chaptermark{Éthique et impartialité}

\epigraph{%
    « L'enquêteur ne cherche pas à avoir raison~;\\
    il cherche ce qui s'est passé. »%
}{Principe du cours}

\niveauChapitre{Fondamental}{Chapitre~\ref{chap:demarche}.}

\begin{boxobjectifs}
\begin{itemize}
    \item Comprendre la neutralité et l'impartialité de l'enquêteur.
    \item Reconnaître les biais qui menacent l'enquête.
    \item Respecter la confidentialité et la proportionnalité.
    \item Gérer les conflits d'intérêts.
    \item Rapporter à charge et à décharge.
\end{itemize}
\end{boxobjectifs}

\bigskip

% ============================================================================
\section*{Question directrice : l'expert doit-il aider à condamner M. NDONGO ?}
% ============================================================================

L'expert n'est ni du côté de l'accusation ni de la défense~: il éclaire. Son
devoir est la vérité, rapportée à charge comme à décharge. Ce chapitre pose
l'éthique de l'enquête.

\medskip
\textbf{Plan du chapitre.}
\begin{enumerate}
    \item Neutralité et impartialité (section~\ref{sec:c03-neutralite}).
    \item Les biais (section~\ref{sec:c03-biais}).
    \item Confidentialité et proportionnalité (section~\ref{sec:c03-confid}).
    \item Conflits d'intérêts (section~\ref{sec:c03-conflits}).
\end{enumerate}

\separateur

% ============================================================================
\section{Neutralité et Impartialité}
\label{sec:c03-neutralite}
% ============================================================================

\begin{boxdef}
\termecle{Impartialité}~: disposition de l'enquêteur à suivre les preuves où
elles mènent, sans favoriser une thèse. \termecle{Neutralité}~: absence
d'intérêt personnel dans l'issue. L'expert rapporte ce qu'il établit, y
compris ce qui disculpe.
\end{boxdef}

% ============================================================================
\section{Les Biais}
\label{sec:c03-biais}
% ============================================================================

\begin{tableaucours}{L{3.6cm}L{8.6cm}}{%
    \tableauHeader{Biais} &
    \tableauHeader{Menace}
}
Biais de confirmation & Ne retenir que ce qui confirme la thèse attendue \\
Pression de l'autorité & Orienter l'analyse vers le résultat voulu \\
Attachement à l'outil & Ne voir que ce que l'outil montre \\
Conclusion hâtive & Interpréter avant d'avoir établi \\
\end{tableaucours}
\vspace{-0.8em}
\begin{center}{\small\itshape Tableau~3.1. Biais qui menacent l'enquête}\end{center}

\begin{boxalert}
\textbf{Violation typique.} Chercher dans le disque « la preuve que NDONGO
est coupable ».

\smallskip
\textbf{Règle~:} on cherche ce qui s'est passé, pas une confirmation. On
formule les hypothèses concurrentes et on les éprouve toutes, à charge et à
décharge.
\end{boxalert}

% ============================================================================
\section{Confidentialité et Proportionnalité}
\label{sec:c03-confid}
% ============================================================================

\begin{boxCRO}
\textbf{Confidentialité et proportionnalité.} Un support saisi contient
souvent des données privées sans rapport avec l'affaire (photos, messages
personnels). Les examiner toutes dégrade la \cro{-}{0}{0}~: on limite
l'analyse au périmètre de la réquisition, et l'on protège le reste. La
proportionnalité est une exigence juridique autant qu'éthique.
\end{boxCRO}

% ============================================================================
\section{Conflits d'Intérêts}
\label{sec:c03-conflits}
% ============================================================================

\begin{boxscenario}{Opération TSANAGA. L'expert se déporte-t-il ?}
\textbf{Constatation.} M.~BIKORO constate qu'il connaît personnellement un
cadre de LiZbiZ impliqué dans l'affaire.

\medskip
\textbf{Analyse} (confiance élevée). Un lien personnel peut entacher
l'impartialité, ou l'apparence d'impartialité. M.~BIKORO le signale au
parquet et à Me~TSALA, qui décident de la suite. La transparence sur les
conflits protège l'enquête.
\end{boxscenario}

\begin{boxmetier}{Expert judiciaire}
\textbf{Servir le tribunal, pas une partie.} L'expert judiciaire doit la
vérité au tribunal, non à qui le mandate. Rapporter un élément à décharge
n'est pas une faiblesse~: c'est le cœur de sa mission.
\end{boxmetier}

\separateur

\begin{boxresume}
\textbf{Ce qu'il faut retenir de ce chapitre~:}
\begin{itemize}
    \item \textbf{Impartialité et neutralité}~: suivre les preuves, sans
          intérêt.
    \item \textbf{Les biais} (confirmation, autorité) menacent l'enquête.
    \item \textbf{Proportionnalité}~: limiter l'analyse au périmètre.
    \item \textbf{Conflits d'intérêts}~: les signaler.
    \item \textbf{À charge et à décharge}~: on rapporte tout.
\end{itemize}
\end{boxresume}

\begin{boxinfo}
\textbf{Points de contrôle de l'élève.} Je sais~: définir impartialité et
neutralité~; nommer les biais~; appliquer la proportionnalité~; gérer un
conflit d'intérêts~; rapporter à décharge.
\end{boxinfo}

\bigskip

\begin{boxexo}[title={\ding{45}\quad Exercice~3.1. Impartialité \niveaubase}]
Qu'est-ce que l'impartialité de l'enquêteur~?
\end{boxexo}

\begin{boxexo}[title={\ding{45}\quad Exercice~3.2. Biais \niveaubase}]
Citez deux biais qui menacent une enquête.
\end{boxexo}

\begin{boxexo}[title={\ding{45}\quad Exercice~3.3. À décharge \niveaubase}]
Pourquoi rapporter un élément qui disculpe le mis en cause~?
\end{boxexo}

\begin{boxexo}[title={\ding{45}\quad Exercice~3.4. Confirmation \niveaumoyen}]
En quoi « chercher la preuve de la culpabilité » est-il un biais~?
\end{boxexo}

\begin{boxexo}[title={\ding{45}\quad Exercice~3.5. Proportionnalité \niveaumoyen}]
Un support contient des photos privées sans lien avec l'affaire. Que faire~?
\end{boxexo}

\begin{boxexo}[title={\ding{45}\quad Exercice~3.6. Conflit \niveaumoyen}]
M.~BIKORO connaît une personne impliquée. Comment réagir~?
\end{boxexo}

\begin{boxexo}[title={\ding{45}\quad Exercice~3.7. Pression \niveaumoyen}]
Comment résister à la pression de l'autorité requérante~?
\end{boxexo}

\begin{boxexo}[title={\ding{45}\quad Exercice~3.8. Apparence \niveaumoyen}]
Pourquoi l'apparence d'impartialité compte-t-elle autant que l'impartialité~?
\end{boxexo}

\begin{boxexo}[title={\ding{45}\quad Exercice~3.9. Charte \niveauexpert}]
Rédigez une charte déontologique de l'expert forensique.
\end{boxexo}

\begin{boxexo}[title={\ding{45}\quad Exercice~3.10. Hypothèses \niveauexpert}]
Formulez, pour TSANAGA, deux hypothèses concurrentes et comment les éprouver.
\end{boxexo}

\bigskip

\begin{boxinfo}
\textbf{Pour aller plus loin.}
\begin{itemize}
    \item ISO/IEC 27037~\cite{iso27037}~; ACPO~\cite{acpo2012}.
\end{itemize}
\end{boxinfo}

\begin{boxinfo}
\textbf{Transition.} L'éthique rejoint le droit~: la Partie~II traite du
cadre juridique et de la recevabilité, à commencer par le droit de la preuve
(Chapitre~\ref{chap:droit}).
\end{boxinfo}
